\documentclass[11pt,a4paper]{article}
\usepackage[T1]{fontenc}
\usepackage{mathptmx}
\usepackage[margin=22mm]{geometry}
\usepackage{amsmath,amssymb,graphicx,booktabs,tabularx,array,float,pdflscape}
\usepackage{microtype}
\usepackage[hidelinks]{hyperref}
\usepackage{url}
\usepackage[font=small,labelfont=bf]{caption}
\newcommand{\figdir}{figures}
\renewcommand{\thefigure}{S\arabic{figure}}
\renewcommand{\thetable}{S\arabic{table}}
\title{Supplementary Information: Historically Constrained Aircraft Design with Language Models}
\author{Ehsan Roohi}
\date{}
\begin{document}
\maketitle
\noindent These source-backed figures and detailed computational tables accompany the main article. Model-authored V0 plates are proposals, not fabricated aircraft; the Wright three-view is a researcher lifting-surface reconstruction. Numerical digits support reproducibility, not measurement precision.
\section*{Supplementary figures}

\begin{landscape}
\begin{figure}[H]\centering
\includegraphics[width=.96\linewidth,height=.75\textheight,keepaspectratio]{\figdir/wright_threeview_reconstruction.png}
\caption{Researcher reconstruction of the 1903 Flyer's lifting surfaces only. Plan, side and front views describe the geometry admitted to the aerodynamic calculation, not the complete historical machine. The moment reference is assumed; the pilot, drive, propellers, rudders, truss and wing-warping deformation are omitted.}\label{fig:wrightthree}
\end{figure}
\end{landscape}

\begin{landscape}
\begin{figure}[H]\centering
\includegraphics[width=.96\linewidth,height=.74\textheight,keepaspectratio]{\figdir/astra_oblique_v0.png}
\caption{Astra's model-authored V0 proposal, shown as an oblique component map. The two supporting surfaces, seated operator, long shaft, aft pusher, rudder and ground gear are proposed features; dashed envelopes and linkages are schematic. The oblique projection does not indicate an off-centre pilot. This is not the V1 analyzed geometry.}\label{fig:astra}
\end{figure}
\end{landscape}

\paragraph{Readable key to the Astra V0 plate.}
The forward surface (F) is a 12.0 m-span, 2.0 m-chord pair of pivoted halves;
their common and differential motion was proposed for pitch and roll,
respectively. The rear supporting surface (A) is 10.0 by 1.6 m and was to
remain fixed in flight. The spanwise pivot bearings (P), vertical rudder (R),
hand/pedal controls (C), and their cable and load paths were not fully
specified. A proposed 72 kg/18 kW motor (M) drives an aft-facing 2.4 m screw
(Q) through a long guarded shaft (S) running below the operator (O); shaft
bearings, coupling and blade geometry remain unknown. Fuel and oil vessels
(V), trim ballast (B), spring vane (I), wheels (W), and front/rear skids (K)
are located only schematically or in a display plane where the true lateral
coordinate is unknown. The dashed brown prism is a truss \emph{envelope}, not
a verified set of wood members and wire attachments. Spars, ribs and fabric
are proposed, while local clearances, structural loads, ground attachment,
pilot reach and control travel remain unresolved. None of these V0 labels
constitutes an installed V12 configuration or a measured performance result.

\begin{landscape}
\begin{figure}[H]\centering
\includegraphics[width=.96\linewidth,height=.80\textheight,keepaspectratio]{\figdir/fable_oblique_v0.png}
\caption{Fable's model-authored V0 biplane proposal. The plate proposes a seated operator, braced upper and lower wings, a rotating horizontal tail, tip roll controls, two counter-rotating pushers and skids. It also retains unresolved mass-centre and tail--fin interference issues; proposed paths are not verified mechanisms.}\label{fig:fable}
\end{figure}
\end{landscape}

\paragraph{Readable key to the Fable V0 plate.}
Two fabric-covered lifting decks are separated by posts and proposed wire
bracing; the upper deck is labelled 11.5 m span by 1.85 m chord. Hinged outer
tip panels of approximately $\pm8^\circ$ were proposed as roll controls.
The seated operator, a central engine group, and two nominal 2.3 m screw
discs occupy the middle airframe. Opposite propeller rotation is stated in
the source but not resolved by the schematic ellipses, and the drive and
control-run paths are not load- or motion-verified. An aft horizontal plane
and fin/rudder were proposed for pitch and yaw, yet the initial rotating-tail
concept conflicts with the fin over some candidate settings. The brown
skids and drawn wheels are only a ground-gear proposal: the source text and
coordinate table disagree on whether there are three or four wheels. Post
stations, arch thickness, pilot restraint, engine support, cable clearance,
and the claimed centre of mass remain assumptions. The V0 picture therefore
documents an initial architecture and its contradictions, not the corrected
V13 mass ledger or an independently trimmed airplane.

\begin{landscape}
\begin{figure}[H]\centering
\includegraphics[width=.96\linewidth,height=.80\textheight,keepaspectratio]{\figdir/opus_oblique_v0.png}
\caption{Opus's model-authored V0 biplane proposal. Two wing decks, roll flaps, seated cockpit, pusher drive, twin booms, aft elevator and rudder, and ground gear are indicated. Control runs and bracing are schematic, and several joint, clearance and ground-contact dimensions are unresolved.}\label{fig:opus}
\end{figure}
\end{landscape}

\paragraph{Readable key to the Opus V0 plate.}
The lower wing is labelled 11.0 m span by 1.8 m chord; an upper wing sits
above it, with outboard balancing flaps for roll. A seated operator and
engine/fuel/radiator group are shown near a pusher screw whose swept disc is
nominally 2.4 m in diameter. The declared shaft stops at $x=2.0$ m while the
disc is at $x=2.1$ m, so the intervening joint is unknown rather than silently
connected. Twin booms and spars carry a tailplane, elevator, fin and rudder;
the red control runs are schematic and do not certify lever reach, hinge
loads or cable clearances. Two ash skids and wheel pairs are proposed, but
the main-skid and tail-skid vertical coordinates conflict with a single
ground datum. The drawing also marks unknown tank/radiator locations,
bracing, throttle location and instrument position. Its original mass labels
and centre-of-mass marker describe the V0 proposal only; they must not be
read as the later source-corrected V12 ledger, installed engine map, or
whole-aircraft trim.

\begin{figure}[H]\centering
\includegraphics[width=\linewidth]{\figdir/wright_response_v1.png}
\caption{Wright reconstruction sensitivity.}\label{fig:wright}
\end{figure}

\begin{figure}[H]\centering
\includegraphics[width=\linewidth]{\figdir/fixed_control_response_v1.png}
\caption{Fixed-control pitch response.}\label{fig:pitch}
\end{figure}

\section*{Supplementary tables}

\begin{table}[H]\centering\small
\caption{Representative interpolated balances, not direct trim validation.}\label{tab:interpolated}
\begin{tabular}{lrrr}\toprule
Quantity & Astra & Fable & Opus\\\midrule
$\delta$, deg &0.0503&-11.2744&-1.8886\\
$V$, m/s &15.438&12.827&16.853\\
$L$, N &3205.180&2836.259&3477.407\\
$T\sin\alpha$, N &31.015&27.283&33.374\\
$W$, N &3236.194&2863.542&3510.781\\
$D=T\cos\alpha$, N &443.534&390.161&477.268\\
$M_{aero}$, N m &-316.128&18.658&207.155\\
$M_T$, N m &316.128&-18.658&-207.155\\
$P_{shaft}$, kW &12.450&10.009&13.406\\
Engine target, kW &18 (unverified)&unknown&13 (unverified)\\\bottomrule
\end{tabular}
\end{table}

\begin{table}[H]\centering\small
\caption{Direct fine-grid lifting-surface loads.}\label{tab:loads}
\begin{tabular}{lrr}\toprule
Surface & Lift, N & Pitch moment about CG, N m\\\midrule
Astra forward &2394.207&1911.535\\
Astra rear &802.779&-2223.350\\
Opus lower deck &1694.552&-107.805\\
Opus upper deck &1901.859&-223.225\\
Opus horizontal tail &-120.841&534.266\\\bottomrule
\end{tabular}
\end{table}

\begin{table}[H]\centering\small
\caption{Direct fine-grid balance at the archived candidate settings.}\label{tab:residual}
\begin{tabular}{lrr}\toprule
Quantity & Astra & Opus\\\midrule
$L+T\sin\alpha$, N &3228.001&3508.944\\
$W$, N &3236.194&3510.781\\
$R_z$, N &-8.194&-1.837\\
$D$, N &442.947&474.920\\
$T\cos\alpha$, N &443.534&477.268\\
$R_x$, N &0.588&2.348\\
$M_{aero}$, N m &-312.049&203.116\\
$M_T$, N m &316.128&-207.155\\
$R_m$, N m &4.079&-4.039\\\bottomrule
\end{tabular}
\end{table}

\begin{table}[H]\centering\small
\caption{Fixed-control pitch diagnostics at the archived candidate settings.}\label{tab:stability}
\begin{tabular}{llrrr}\toprule
Design & Mesh & $C_{L_\alpha}$/deg & $C_{m_\alpha}$/deg & $SM_{L,eff}$, \% $c_{ref}$\\\midrule
Astra &8/40&0.061280&0.005980&-9.758\\
Astra &12/80&0.061180&0.006080&-9.938\\
Opus &8/40&0.065380&-0.003280&5.017\\
Opus &12/80&0.065440&-0.003360&5.134\\\bottomrule
\end{tabular}
\end{table}

\begin{table}[H]\centering\footnotesize
\caption{V1 architectural comparison and remaining evidence.}\label{tab:architecture}
\begin{tabularx}{\linewidth}{l*{4}{>{\raggedright\arraybackslash}X}}\toprule
System & Wright 1903 & Astra & Fable & Opus\\\midrule
Lifting layout & Main biplane, forward elevator & Tandem lifting surfaces & Biplane, aft tail & Biplane, aft tail\\
Pilot/control & Prone; mechanical pitch, warping and yaw controls & Seated; collective/differential forward incidence & Seated; movable tips and rotating tail & Seated; flaps, elevator and rudder\\
Propulsion & Twin counter-rotating pushers; historical operation & Pusher with long shaft; installation unverified & Twin counter-rotating pushers; available power unresolved & Pusher; conditional power margin inadequate at selected efficiency\\
Structure & Historical built artifact; later replicas require separate identity & Truss/wing-root and bearing paths incomplete & Lattice/member and joint definition incomplete & Spar/strut/wire definition incomplete\\
Evidence & Documented flights, later tests and reconstructions & Numerical conditions only & Mechanical conflict in sampled trims & Numerical conditions only\\\bottomrule
\end{tabularx}
\end{table}

\begin{table}[H]\centering\small
\caption{Selected V2 longitudinal screening states. Astra and Opus use different model-proposed operating points; Fable remains blocked before independent trim. Engine demands use assumed propeller and drive efficiencies and do not establish installed propulsion or flight.}\label{tab:v2}
\begin{tabular}{lrrl}\toprule
Quantity & Astra V2 & Opus V2 & Fable V2\\\midrule
Mass, kg & 320 & 360 & 295 (ledger)\\
Speed, m/s & 18.0 & 15.0 & --\\
$\alpha$, deg & $-0.103$ & 6.347 & --\\
Pitch control, deg & $-0.330$ & $-3.168$ & --\\
$R_z$, N & $-0.063$ & $-0.025$ & --\\
$R_m$, N m & $+0.065$ & $+0.0003$ & --\\
$SM_{L,eff}$, \% $c_{ref}$ & 19.47 & 7.76 & --\\
Useful power, kW & 4.42 & 6.65 & --\\
Engine demand, kW, assumed & 8.46 & 11.67 & --\\\bottomrule
\end{tabular}
\end{table}

\begin{table}[H]\centering\footnotesize
\caption{Version inheritance and evaluation status. V0 is the initial concept; V1 is the frozen clarification used for the main screen; V2 received modern case-specific feedback. All stated improvements are hypotheses until independently checked.}\label{tab:stage}
\begin{tabularx}{\linewidth}{lXXX}\toprule
Case & V0 concept & V1 inherited or clarified & V2 changed; independent status\\\midrule
Astra & Tandem lift, seated pilot, long aft-pusher shaft & Surface envelopes retained; 330 kg leaf ledger and analysis CG replace pictorial ballast position & Fixed wing/aft tail, tractor and short shaft; conditional longitudinal retrim only\\
Fable & Biplane, twin pushers, continuous aft tail; contradictory CG & 292 kg leaf ledger, shifted deck stations, continuous tail; tail/fin collision at sampled trim & Split tail, relocated fin and revised drive; geometry and propulsion gates still open\\
Opus & Biplane, single pusher, aft controls & 358 kg leaf ledger and defined 38 m$^2$ reference; conditional V1 trim & 360 kg ledger, longer skids, outboard booms, revised rudder and slower proposed state; conditional longitudinal retrim only\\\bottomrule
\end{tabularx}
\end{table}

\begin{table}[H]\centering\footnotesize
\caption{Flight-dynamics derivative screen at the frozen V2 trim states. Angular slopes in the first row are per degree; all reported rate derivatives use the nondimensional rates defined in the main article. Coarse/fine variation is numerical sensitivity, not validation. The geometry omits body/fin and installed propulsion; missing derivatives are not zero.}\label{tab:derivativegate}
\begin{tabularx}{\linewidth}{lXXX}\toprule
Quantity & Astra V2 & Fable V2 & Opus V2\\\midrule
$C_{L_\alpha},C_{m_\alpha}$ & 0.07292, $-0.01420$ per degree & No independent retrim & 0.06444, $-0.00500$ per degree\\
$C_{m_q}$ coarse/fine & $-4.6856/-4.7059$ & Not evaluated & $-6.6103/-6.6458$\\
$C_{L_q}$ coarse/fine & $6.2869/6.3022$ & Not evaluated & $5.5957/5.6109$\\
$C_{l_p}$ fine & $-0.3834$ & Not evaluated & $-0.3746$\\
$C_{n_r}$ fine, partial geometry & $-0.0076$ & Not evaluated & $-0.0124$\\
$C_{n_\beta}$ fine, partial geometry & approximately 0 & Not evaluated & $-0.00084$\\
$C_{m_{\dot\alpha}},C_{L_{\dot\alpha}}$ & Downwash-delay model absent & Downwash-delay model absent & Downwash-delay model absent\\
Installed directional/control derivatives & Body/fin and full control model absent & Geometry gate open & Body/fin and full control model absent\\
Full inertia and dynamic eigenmodes & Intrinsic component inertia absent & Intrinsic component inertia absent & Intrinsic component inertia absent\\\bottomrule
\end{tabularx}
\end{table}

\begin{table}[H]\centering\footnotesize
\caption{Explicitly reduced $\alpha$--pitch-rate sign screen for the independently trimmed V2 lifting-surface models. $U$ is trim speed; $Z_\alpha$, $Z_q$, $M_\alpha$, $M_q$ and the two-state matrix $A_{\alpha q}$ are defined in the main-text reduced-system equation. Products with $I_{yy}$ permit the sign check without assigning the unknown complete pitch inertia. A negative trace and positive determinant hold for any finite $I_{yy}>0$ in this \emph{omission-based} two-state matrix only. This is not a computed aircraft eigenmode, full dynamic-stability result, or flight clearance.}\label{tab:reduceddynamic}
\begin{tabular}{lrr}\toprule
Diagnostic & Astra V2 & Opus V2\\\midrule
$Z_\alpha/U$, s$^{-1}$ & $-4.3191$ & $-3.5807$\\
$1+Z_q/U$, dimensionless & $0.5476$ & $0.6735$\\
$I_{yy}M_\alpha$, N m rad$^{-1}$ & $-12100.5$ & $-2699.4$\\
$I_{yy}M_q$, N m s & $-4864.0$ & $-3758.8$\\
$I_{yy}\det(A_{\alpha q})$, kg m$^2$ s$^{-2}$ & $27634.7$ & $15277.0$\\
Full longitudinal/lateral eigenmodes & Not assessed & Not assessed\\\bottomrule
\end{tabular}
\end{table}
\begin{table}[H]\centering\footnotesize
\caption{Initial versus revised pitch-derivative evidence at each version's own conditional trim speed. V1 rate derivatives come from one retained AVL panelization; V2 values use the finer of two panelizations. All are rigid lifting-surface, quasi-steady outputs, not installed-aircraft or measured coefficients. The reduced two-state sign check omits unsteady, axial-speed, propulsive and control couplings. Therefore an affirmative result in its row is not a pass of the full flight-dynamics requirements. Fable has no mechanically admissible V1 trim and no independently reconciled V2 integrated geometry.}\label{tab:v1v2dynamic}
\begin{tabular}{lrrrr}\toprule
Quantity & Astra V1 & Astra V2 & Opus V1 & Opus V2\\\midrule
Trim speed, m s$^{-1}$ & 15.461 & 18.000 & 16.859 & 15.000\\
$C_{m_\alpha}$, rad$^{-1}$ & $+0.3483$ & $-0.8130$ & $-0.1921$ & $-0.2864$\\
$C_{m_q}$ & $-5.7865$ & $-4.7059$ & $-6.5909$ & $-6.6458$\\
$C_{L_q}$ & $4.7580$ & $6.3022$ & $5.6419$ & $5.6109$\\
Local static restoring slope & No & Yes & Yes & Yes\\
Reduced two-state sign check & Yes & Yes & Yes & Yes\\
Full dynamic acceptance & Not assessed & Not assessed & Not assessed & Not assessed\\\bottomrule
\end{tabular}
\end{table}
\paragraph{Post-analysis integrated V3 returns.}
The case-specific V3 prompts requested a new complete-aircraft proposal and a
flight-dynamics evidence plan, not an assurance of flight. Astra completed on
the first bounded call. Fable and Opus were cut off at the caller's 20,000-token
output limit; the incomplete responses were retained, and independent
full-prompt retries at a 48,000-token limit returned complete texts. The
structured mass ledgers below were re-summed from their individual entries
with $M=\sum_i m_i$ and $\mathbf{x}_{CG}=\sum_i m_i\mathbf{x}_i/M$.
The native longitudinal axes differ, so $x_{CG}$ values must not be compared
without frame transformation. The returned masses and coordinates are
model assumptions, not measurements. All request and response hashes and the
arithmetic auditor are in the public archive. No V3 numerical result below
replaces the independently computed V1/V2 results in the main analysis.

\begin{table}[H]\centering\footnotesize
\caption{Completed V3 model returns and limited independent arithmetic gate. CG coordinates are in each model's native frame, in metres. ``Force/moment closure'' means arithmetic consistency of the model's stated loads only; it is not an independently predicted aerodynamic trim. No candidate has an independently established installed derivative set, complete longitudinal/lateral roots, or flight acceptance.}\label{tab:v3returns}
\begin{tabularx}{\linewidth}{lXXX}\toprule
Return & Recomputed mass and CG & Nominal force/moment evidence & Full dynamic verdict\\\midrule
Astra V3-I & 333.0 kg; $(-0.150889,0,0.218445)$ & Actual powered/glide trim unknown & Roots unknown; not assessed\\
Fable V3 & 309.5 kg; $(1.556656,0,0.765670)$ & $3245-209-3036=0$ N; $-1329+436-14+907=0$ N m, using assumed loads & Reduced estimates only; not assessed\\
Opus V3 & 361.3 kg; $(0.622676,0,0.460036)$ & Trim left to evaluator by model & Lateral roots unknown; not assessed\\\bottomrule
\end{tabularx}
\end{table}
\end{document}
